\section{Boxed Environments}

This template includes a beautiful boxed environment for problems and a cleanly separated environment for solutions. It also includes environments for definitions, theorems, and notes.

\subsection{Problem Environment}
Use the \texttt{problem} environment to state your questions. You can provide an optional title to the problem. 

\begin{problem}[1: Signal Analysis]
Calculate the energy and power of the given signal $x(t) = e^{-2t}u(t)$.
\end{problem}

\subsection{Solution Environment}
Following the problem, you can use the \texttt{solution} environment to provide your answer. It will automatically style the section and add a QED marker at the end.

\begin{solution}
The energy of the signal is given by:
\begin{align*}
E &= \int_{-\infty}^{\infty} |x(t)|^2 \,dt \\
  &= \int_{0}^{\infty} e^{-4t} \,dt \\
  &= \left[ \frac{e^{-4t}}{-4} \right]_0^\infty \\
  &= \frac{1}{4} \text{ Joules}
\end{align*}
Since the energy is finite ($E < \infty$), the power is $P = 0$.
\end{solution}

\subsection{Other Boxed Environments}
You can use the \texttt{definition}, \texttt{theorem}, and \texttt{note} environments to highlight important concepts.

\begin{definition}[Energy Signal]
A signal is an energy signal if its total energy satisfies $0 < E < \infty$.
\end{definition}

\begin{theorem}[Parseval's Theorem]
The total energy computed in the time domain is equal to the total energy computed in the frequency domain.
\end{theorem}

\begin{note}[Important]
Power signals have infinite energy, while energy signals have zero average power.
\end{note}

\newpage